\documentclass[final]{beamer}
\usepackage[orientation=portrait,size=a1,scale=1.1]{beamerposter}
\usepackage[brazilian]{babel}
\usepackage{amsmath}
\usetheme{default}
\title{Modelagem computacional reprodutível de estruturas moleculares de carbono}
\author{Aloisio Magalhães}
\institute{Projeto de físico-química computacional}
\begin{document}
\begin{frame}[t]
\maketitle
\begin{columns}[t]
\column{.48\textwidth}
\begin{block}{Pergunta}
Como um pipeline reprodutível compara energia e propriedades eletrônicas de modelos pequenos de carbono sob diferentes aproximações de estrutura eletrônica?
\end{block}
\begin{block}{Método}
Python + PySCF; base STO-3G; HF, PBE e B3LYP; sistemas H2O, C2, C4 linear, anel C6 e C4 tetraédrico. O GitHub Actions registra código, ambiente, logs e artefatos.
\end{block}
\begin{block}{Observáveis}
Energia total, convergência, HOMO, LUMO, gap, cargas de Mulliken e geometria. A saída principal é JSON; a síntese é CSV e gráfico.
\end{block}
\column{.48\textwidth}
\begin{block}{Interpretação}
Os modelos servem como benchmark e treinamento de reprodutibilidade. C4 linear e C6 são aproximações moleculares de estruturas sp2; C4 tetraédrico é uma aproximação de coordenação sp3. Nenhum substitui um cálculo periódico.
\end{block}
\begin{block}{Limitações}
Não há periodicidade, crescimento cristalino, amostra de cinza, radioisótopo ou bateria. Falha de convergência é registrada. Uma previsão computacional não é prova experimental.
\end{block}
\begin{block}{Reprodução}
Repositório: \url{https://github.com/AloisioMagalhaes/fisico-quimica-computacional-piaui}

Workflow: \url{https://github.com/AloisioMagalhaes/fisico-quimica-computacional-piaui/actions/workflows/quantum.yml}
\end{block}
\end{columns}
\end{frame}
\end{document}
